\documentclass[10pt,a4paper]{amsart}
\usepackage[margin=19mm]{geometry}
\usepackage{amsmath,amssymb,mathtools}
\usepackage[scr=rsfs,cal=euler]{mathalfa}
\usepackage{xfrac}
\usepackage[unicode,hidelinks,bookmarks=false]{hyperref}
\newcommand{\ie}{i.e.\ }
\newcommand{\cf}{cf.\ }
\newcommand{\resp}{respectively\ }
\newcommand{\Subsection}{\subsection}
\providecommand{\nobreakdash}{\nobreak\hbox{-}}
\setlength{\emergencystretch}{2em}
\multlinegap0pt
\usepackage{bm}
\usepackage{bbm}
\usepackage{dsfont}
\usepackage{xspace}
\usepackage{cancel}
\usepackage{mathtools}
\usepackage{amsthm}
\makeatletter
\@ifundefined{claim}{\newtheorem*{claim}{Claim}}{}
\makeatother
\def\Tr{\text{Tr }}
\title[Optimization of maximal quantum f-divergences between unitar]{Optimization of maximal quantum f-divergences between unitary orbits}
\author{Hoang Minh Nguyen, Hoang An Nguyen, Cong Trinh Le}
\date{Source: arXiv:2601.08268v2 (CC BY 4.0)}
\begin{document}
\maketitle

\noindent\emph{Source and licence.} Optimization of maximal quantum f-divergences between unitary orbits. Version arXiv:2601.08268v2 (CC BY 4.0), distributed under the Creative Commons Attribution 4.0 International licence, as recorded in the arXiv licence metadata for this exact version and shown on the abstract page https://arxiv.org/abs/2601.08268v2. Full rights notice: licenses/arxiv-2601.08268-CC-BY-4.0.txt.\par\smallskip

\noindent\textit{Editorial scope.} Counted unit: the Claim whose source label is \texttt{claimrho} in the article (source0.tex lines 794--801, source label \texttt{claimrho}), delivered together with the article's own complete original proof of it (source0.tex lines 802--993, one \texttt{proof} environment of 192 lines). No mathematics is written, repaired or re-derived: statement and proof are the article's own text line for line. Required context retained uncounted: claimA (lines 597--604). Every definition, display and label of those lines is retained unchanged. Omitted from this package and not redistributed: 1--596, 605--793, 994--1142. Nothing inside a retained excerpt is omitted or altered: every retained body is delivered exactly as the article's lines are written, so no omission receipt of any kind accompanies this package. Citations: the retained excerpts carry no citation command at all, so no bibliography entry of the article is transcribed and no reference is redistributed. Reference adaptation: none was needed and none was made; every reference inside the delivered unit resolves inside it, so no unresolved reference or citation is produced. Printed slips kept literal: no printed word, symbol, hypothesis or number of the counted statement or of its proof was altered. Layout and class: the article's own class is \texttt{amsart} with packages including amsmath, amssymb, amsthm, graphicx, xypic, vietnam, hyperref, chngcntr; the wrapper is the lane's pinned \texttt{amsart} editorial wrapper at 10pt on A4, which restates the theorem-like environments the retained text needs and adds the article's own macro definitions that the retained text needs (\texttt{\textbackslash Tr}), with extra wrapper packages (bm, bbm, dsfont, xspace, cancel, mathtools). The delivered numbering is the unit's own. Assets: no retained span contains a figure environment, a graphics-inclusion command, a PDF-page inclusion, a table or a TikZ picture; no image file is copied into this package, the package declares no asset dependency and no image file is redistributed with it. Licence: version arXiv:2601.08268v2 of this article is distributed under the Creative Commons Attribution 4.0 International licence, as recorded in the arXiv licence metadata for this exact version and shown on the abstract page https://arxiv.org/abs/2601.08268v2. A full rights notice is delivered at licenses/arxiv-2601.08268-CC-BY-4.0.txt. Characters: every retained excerpt is pure ASCII, so the delivered engine, XeLaTeX with its default Latin Modern text font, reports no missing character.
\begin{claim}\label{claimA} 
\begin{align*}
\max_{U\in\mathbb U_n} \operatorname{Tr}\dfrac{(U^*\sigma^{-1}U)\rho^2}{(U^*\sigma^{-1/2}U)\rho(U^*\sigma^{-1/2}U)+sI}&=\operatorname{Tr}\dfrac{\lambda^\uparrow(\sigma)^{-1}\lambda^\downarrow(\rho)^2}{\lambda^\uparrow(\sigma)^{-1}\lambda^\downarrow(\rho) +s I};\\
\arg\max_{U\in \mathbb U_n}\operatorname{Tr}\dfrac{(U^*\sigma^{-1}U)\rho^2}{(U^*\sigma^{-1/2}U)\rho(U^*\sigma^{-1/2}U)+sI}&=W^\uparrow V^{\downarrow*};\\
\min_{U\in\mathbb U_n}\operatorname{Tr}\dfrac{(U^*\sigma^{-1}U)\rho^2}{(U^*\sigma^{-1/2}U)\rho(U^*\sigma^{-1/2}U)+sI}&= \operatorname{Tr}\dfrac{\lambda^\downarrow(\sigma)^{-1}\lambda^\downarrow(\rho)^2}{\lambda^\downarrow(\sigma)^{-1}\lambda^\downarrow(\rho) +s I};\\
\arg\min_{U\in \mathbb U_n} \operatorname{Tr}\dfrac{(U^*\sigma^{-1}U)\rho^2}{(U^*\sigma^{-1/2}U)\rho(U^*\sigma^{-1/2}U)+sI}&=W^\downarrow V^{\downarrow*}.
\end{align*}
\end{claim}

\begin{claim}\label{claimrho}
\begin{align*}
\max_{U\in\mathbb U_n} -2 \operatorname{Tr}\dfrac{\rho}{(U^*\sigma^{-1/2}U)\rho(U^*\sigma^{-1/2}U)+sI}&=-2\operatorname{Tr}\dfrac{\lambda^\downarrow(\rho)}{\lambda^\uparrow(\sigma)^{-1}\lambda^\downarrow(\rho)+sI};\\
\arg\max_{U\in\mathbb U_n} -2 \operatorname{Tr}\dfrac{\rho}{(U^*\sigma^{-1/2}U)\rho(U^*\sigma^{-1/2}U)+sI}&=W^\uparrow V^{\downarrow*};\\
\min_{U\in\mathbb U_n} -2 \operatorname{Tr}\dfrac{\rho}{(U^*\sigma^{-1/2}U)\rho(U^*\sigma^{-1/2}U)+sI}&=-2\operatorname{Tr}\dfrac{\lambda^\downarrow(\rho)}{\lambda^\downarrow(\sigma)^{-1}\lambda^\downarrow(\rho)+sI};\\
\arg\min_{U\in\mathbb U_n} -2 \operatorname{Tr}\dfrac{\rho}{(U^*\sigma^{-1/2}U)\rho(U^*\sigma^{-1/2}U)+sI}&=W^\downarrow V^{\downarrow*}.
\end{align*}
\end{claim}

\begin{proof} 
For $U\in\mathbb U_n$, define
$$
\Phi(U)
:=
-2\,\Tr\!\left[
\rho\Big((U^*\sigma^{-1/2}U)\rho(U^*\sigma^{-1/2}U)+sI\Big)^{-1}
\right].
$$
Set
$$ B:=U^*\sigma^{-1/2}U,\quad A:=U^*\sigma^{-1}U=B^2,$$
so that
\[
\Phi(U)=-2\,\Tr\!\left[\rho\,(A^{1/2}\rho A^{1/2}+sI)^{-1}\right].
\]
Since $U\mapsto A$ ranges over the unitary orbit of $\sigma^{-1}$, the problem is equivalent to optimizing over $A$ belongs to the unitary orbit $\mathbb U_{\sigma^{-1}}$ of $\sigma^{-1}$.

Let
$$
J(A):=\Tr\!\big[\rho\,(A^{1/2}\rho A^{1/2}+sI)^{-1}\big],
\qquad A\in\mathbb U_{\sigma^{-1}},
$$
and define
$$
X:=A^{1/2}\rho A^{1/2},\quad
Y:=(X+sI)^{-1}.
$$
Note that $X$ and $Y$ are Hermitian, with $Y\succ 0$.

Let $K$ be an arbitrary skew-Hermitian matrix ($K^*=-K$), and consider the curve
$$
A(t)=e^{-tK}Ae^{tK},
$$
which lies entirely in the unitary orbit of $A$.
By functional calculus,
$$
A(t)^{1/2}=e^{-tK}A^{1/2}e^{tK},
$$
and hence
$$
\dot A^{1/2}(0)=[A^{1/2},K]=:M,
\qquad M^*=-M.
$$
Since $X(t)=A(t)^{1/2}\rho A(t)^{1/2}$,
$$
\dot X(0)=\dot A^{1/2}(0)\rho A^{1/2}+A^{1/2}\rho\dot A^{1/2}(0)
= M\rho A^{1/2}+A^{1/2}\rho M.
$$
Using the identity
\(
\frac{d}{dt}(Z(t)^{-1})=-Z(t)^{-1}\dot Z(t)Z(t)^{-1},
\)
we obtain
\[
\dot Y(0)=-Y\,\dot X(0)\,Y.
\]
Differentiating $J(A(t))=\Tr(\rho Y(t))$ yields
$$
\dot J(0)
=\Tr(\rho\dot Y(0))
=-\Tr\big(\rho Y\,\dot X(0)\,Y\big).
$$
Substituting the expression for $\dot X(0)$ and using cyclicity of the trace,
\begin{align*}
\dot J(0)
&=-\Tr\big(\rho Y (M\rho A^{1/2}+A^{1/2}\rho M)Y\big)\\
&=-\Tr\Big(M\big(\rho A^{1/2}Y\rho Y+Y\rho Y A^{1/2}\rho\big)\Big).
\end{align*}
Now we define
$$
T:=\rho A^{1/2}Y\rho Y+Y\rho Y A^{1/2}\rho.
$$
Since $A^{1/2}$, $\rho$, and $Y$ are Hermitian and
\[
(\rho A^{1/2}Y\rho Y)^*=Y\rho Y A^{1/2}\rho,
\]
it follows that $T$ is Hermitian. With this definition,
\[
\dot J(0)=-\Tr(MT),
\qquad M=[A^{1/2},K].
\]
Using the trace identity
\[
\Tr([A^{1/2},K]\,T)=\Tr\big(K\,[T,A^{1/2}]\big),
\]
we may rewrite
\[
\dot J(0)=-\Tr\big(K\,[T,A^{1/2}]\big).
\]
At a local extremum of $J$ on the unitary orbit, $\dot J(0)=0$ for all
skew-Hermitian $K$. Since $[T,A^{1/2}]$ is itself skew-Hermitian, this implies
\[
[T,A^{1/2}]=0,
\quad\text{equivalently}\quad
[A^{1/2},T]=0.
\]

We now show that the commutation relation $[A^{1/2},T]=0$ implies $[A,\rho]=0$.
Since $Y$ is obtained from $X$ by functional calculus, $X$ and $Y$ have the same
spectral projections and therefore the same commutant
\[
\{X\}'=\{Y\}'.
\]

Using $X=A^{1/2}\rho A^{1/2}$, the matrix $T$ can be rewritten as
\[
T
=
\rho A^{1/2}Y\rho Y+Y\rho Y A^{1/2}\rho
=
A^{-1/2}X A^{-1} X A^{-1/2}Y
-
X\,Y\,A^{-1/2}X A^{-1/2}Y.
\]
Thus $T$ belongs to the $\ast$-algebra generated by $X$, $Y$, and $A^{\pm 1/2}$.
The relation $[A^{1/2},T]=0$ implies that $A^{1/2}$ commutes with all spectral
projections of $T$. Since $Y$ and $X$ share the same spectral projections and
$T$ contains $Y$ as a nontrivial factor, it follows that $A^{1/2}$ must commute
with the spectral projections of $X$, and hence
\[
[A^{1/2},X]=0.
\]

Consequently,
\[
0=[A^{1/2},X]
=[A^{1/2},A^{1/2}\rho A^{1/2}]
=A^{1/2}[A^{1/2},\rho]A^{1/2}.
\]
Because $A^{1/2}$ is invertible, this implies $[A^{1/2},\rho]=0$, and therefore
\[
[A,\rho]=[(A^{1/2})^2,\rho]=0.
\]
Therefore, at the maximum and minimum, we may assume $[A, \rho] = 0$.


Let
\[
\rho=\operatorname{diag}(r_1,\dots,r_n),
\qquad
r_1\ge\cdots\ge r_n>0,
\]
and
\[
A=\operatorname{diag}(a_1,\dots,a_n),
\]
where $(a_1,\dots,a_n)$ is a permutation of the eigenvalues
$(d_1,\dots,d_n)$ of $\sigma^{-1}$.
Then
\[
\Phi(U)
=
-2\sum_{i=1}^n \frac{r_i}{a_i r_i+s}.
\]
Thus the optimization reduces to a permutation problem
\[
\Phi_\pi
=
-2\sum_{i=1}^n f(d_{\pi(i)},r_i),
\qquad
f(d,r):=\frac{r}{dr+s}.
\]

A direct computation yields
\[
\frac{\partial^2}{\partial r\,\partial d}f(d,r)
=
-\frac{2rs}{(dr+s)^3}<0,
\qquad d,r,s>0.
\]
Hence $f$ is strictly submodular.  Similar to the argument presented in the proof of Claim \ref{claimA}, by the rearrangement principle for submodular functions, $\sum_i f(d_{\pi(i)},r_i)$ is minimized when
$(d_{\pi(i)})$ and $(r_i)$ are ordered in the same sense, and maximized when they
are ordered in opposite senses.

Since $(r_i)=\lambda^\downarrow(\rho)$, the minimum of $\sum_i f(d_{\pi(i)},r_i)$ is attained for
$d_{\pi(i)}=\lambda^\uparrow(\sigma)^{-1}_i$, and the maximum for
$d_{\pi(i)}=\lambda^\downarrow(\sigma)^{-1}_i$.
Multiplying by $-2$ gives the stated formulas for
$\max\Phi(U)$ and $\min\Phi(U)$.

Let $V^\downarrow$ diagonalize $\rho$ with eigenvalues in decreasing order,
and let $W^\uparrow$ (resp.\ $W^\downarrow$) diagonalize $\sigma$ with eigenvalues
in increasing (resp.\ decreasing) order.
Then
\[
U_{\max}=W^\uparrow V^{\downarrow *},
\qquad
U_{\min}=W^\downarrow V^{\downarrow *}
\]
produce the required eigenvalue matchings and hence attain the maximum and minimum,
respectively.
\end{proof}

\end{document}
